\chapter{Descripción Funcional del Payload}
\label{ch:funcional}

Este capítulo presenta la descripción funcional del \textbf{CMOS Microscopy
Payload (CMOSM)}, incluyendo la arquitectura general del sistema, la
descomposición en bloques funcionales y el flujo de datos desde la adquisición
hasta el almacenamiento. El objetivo es describir \textbf{cómo está organizado
internamente el sistema} y cómo interactúan sus componentes de hardware y
software durante la operación nominal y experimental.

% -------------------------------------------------------------------
\section{Visión general de la arquitectura}

El CMOSM está concebido como un sistema compacto de adquisición de imágenes,
basado en un sensor CMOS comercial y un sistema de iluminación programable,
controlados por una computadora de carga útil (Raspberry Pi). La arquitectura
funcional se divide en cuatro bloques principales:

\begin{itemize}
    \item Bloque de adquisición óptica,
    \item Bloque de iluminación,
    \item Bloque de control y procesamiento,
    \item Bloque de almacenamiento y datos.
\end{itemize}

Esta arquitectura permite separar claramente las funciones físicas de captura
de imagen de las funciones de control, procesamiento y análisis, facilitando la
verificación y la trazabilidad del sistema.

\begin{figure}[H]
    \centering
    \includegraphics[width=0.85\textwidth]{figures/fig_architectura_general_cmosm.png}
    \caption{Arquitectura funcional del CMOS Microscopy Payload, mostrando los cuatro bloques principales y el flujo de datos.}
    \label{fig:arquitectura_cmosm}
\end{figure}

% -------------------------------------------------------------------
\section{Descomposición funcional del sistema}

\subsection{Bloque de adquisición óptica}

El bloque de adquisición óptica está compuesto por el sensor CMOS OV5647,
operando en configuración de microscopía sin lentes (lensless) y en contacto
cercano con la muestra. La formación de imagen se basa en la interacción directa
de la luz transmitida o dispersada por la muestra con el plano del sensor.

Las funciones principales de este bloque son:
\begin{itemize}
    \item Conversión de fotones incidentes en señal eléctrica.
    \item Integración temporal de carga controlada mediante el parámetro
    \texttt{ExposureTime}.
    \item Digitalización de la señal y entrega de frames al sistema de control.
\end{itemize}

El sensor opera bajo control manual de exposición y ganancia, deshabilitando
mecanismos automáticos para garantizar consistencia entre capturas.

% -------------------------------------------------------------------
\subsection{Bloque de iluminación}

El bloque de iluminación utiliza una pantalla OLED de 1.3 pulgadas controlada
vía bus I2C como fuente de iluminación plana y programable. Este bloque permite
la generación de patrones simples de iluminación, tanto uniformes como
estructurados.

Sus funciones incluyen:
\begin{itemize}
    \item Generación de patrones de iluminación estáticos.
    \item Control del nivel relativo de intensidad luminosa.
    \item Conmutación rápida de patrones para modos experimentales.
\end{itemize}

En operación nominal, el sistema emplea patrones de iluminación uniformes para
optimizar el contraste de la muestra. La iluminación estructurada se utiliza
exclusivamente en modos experimentales.

% -------------------------------------------------------------------
\subsection{Bloque de control y procesamiento}

El bloque de control y procesamiento está implementado en la Raspberry Pi,
ejecutando el software principal del payload. Este bloque coordina la operación
de todos los subsistemas y gestiona la lógica de estados descrita en el
Capítulo~\ref{ch:modos-operacion}.

Las funciones principales incluyen:
\begin{itemize}
    \item Configuración del sensor CMOS mediante la librería Picamera2.
    \item Control del sistema de iluminación OLED.
    \item Gestión de los modos operativos (Init, Preview, Capture, Safe).
    \item Ejecución de calibraciones básicas (dark frame y flat field).
    \item Cálculo de métricas rápidas de calidad de imagen.
\end{itemize}

Este bloque actúa como interfaz entre el hardware de captura y el bloque de
almacenamiento de datos.

% -------------------------------------------------------------------
\subsection{Bloque de almacenamiento y datos}

El bloque de almacenamiento se encarga de la persistencia de la información
generada por el payload. Incluye tanto las imágenes adquiridas como los
metadatos asociados a cada captura.

Las funciones de este bloque son:
\begin{itemize}
    \item Almacenamiento de imágenes en formatos sin compresión (PNG/TIFF).
    \item Generación de archivos de metadatos en formato estructurado (JSON).
    \item Organización de archivos siguiendo una convención temporal estándar.
\end{itemize}

Este diseño asegura trazabilidad completa de cada experimento y facilita el
análisis posterior en tierra.

% -------------------------------------------------------------------
\section{Flujo funcional de datos}

El flujo funcional del CMOSM durante una captura nominal puede resumirse como:

\begin{enumerate}
    \item La iluminación OLED emite el patrón seleccionado.
    \item La luz interactúa con la muestra y alcanza el sensor CMOS.
    \item El sensor integra la señal durante el tiempo de exposición definido.
    \item El frame digital es transferido al bloque de control.
    \item Se aplican correcciones básicas (dark y flat).
    \item Se calculan métricas de calidad de imagen.
    \item La imagen y los metadatos son almacenados.
\end{enumerate}

Este flujo se mantiene constante tanto para capturas nominales como
experimentales, variando únicamente la secuencia de iluminación y el número de
frames adquiridos.

% -------------------------------------------------------------------
\section{Interfaz funcional con el sistema OBDH}

El CMOSM se comunica con el sistema de OBDH del CubeSat INTISAT a nivel funcional
mediante comandos de control y estados de operación. En la fase actual del
proyecto, esta interfaz se limita a:

\begin{itemize}
    \item Comandos de inicio y detención de captura.
    \item Señales de estado operativo (Idle, Capturing, Safe).
    \item Sincronización temporal básica.
\end{itemize}

El manejo detallado de telemetría y enlaces de datos se considera fuera del
alcance del presente CDR y será abordado en fases posteriores del proyecto.

% -------------------------------------------------------------------
\section{Consideraciones de diseño funcional}

Desde el punto de vista funcional, el diseño del CMOSM prioriza:

\begin{itemize}
    \item Simplicidad operativa y reproducibilidad experimental.
    \item Separación clara entre funciones nominales y experimentales.
    \item Control manual de parámetros críticos del sensor.
    \item Trazabilidad completa de los datos generados.
\end{itemize}

Estas decisiones permiten validar el sistema como demostrador tecnológico de
microscopía sin lentes y facilitan futuras extensiones del diseño.
